The Cram\'er distance between two fixed CDFs is a proper divergence: it is non-negative and equals zero if and only if the two distributions are identical \citep{szekely2013energy,thorarinsdottir2013using}. This is a weaker guarantee than the classical propriety of \S2---it says that the best forecast \emph{for a given observation} is the observation itself, but does not directly address what happens when the observation is stochastic.

The substantive question is whether the pipeline---constructing $F_\text{obs}$ from a noisy draw $y$ and scoring $F_\text{pred}$ against it---rewards forecasts that get the location right. I verify this numerically: drawing $N = 500$ observations from a known distribution, constructing $F_\text{obs}$ for each draw, and sweeping a forecast parameter $\mu$ across a grid. The argmin falls at the true $\mu$ to within grid resolution for all tested values across three distributional families (log-normal, Gaussian, count-like); a refined sweep at $\Delta\mu = 0.005$ with $N = 2{,}000$ draws confirms residual bias $\leq 0.01$ (Appendix~\ref{sec:appendix_consistency}). One limiting case is analytically tractable: as $\sigma_\text{rel} \to 0$, each $F_\text{obs}(y)$ converges pointwise to the step function at $y$, so the Cram\'er distance reduces to classical CRPS by dominated convergence, and location-consistency follows from strict propriety \citep{gneiting2007strictly}. Underestimation of observational uncertainty therefore degrades gracefully to the classical regime. A formal proof of location-consistency for the general case via the energy-distance framework is beyond the scope of this paper.

Scale-consistency does \emph{not} hold: the minimising spread absorbs the observation noise, just as a blurred photograph combines camera-shake and lens-blur into a single effective width. For conflict fatality data, this effect is more consequential than in meteorology or hydrology: the RVI model of \citet{vesco2026uncertainty} implies observational uncertainty of $\sigma_\text{rel} \approx 0.2$--$0.7$ in log-space, which can equal or exceed forecast-model spread. The metric rewards forecasts that match the observable data distribution, not the latent truth---the same property classical CRPS has in degenerate form. This does not affect the ranking reversals demonstrated in \S6, which depend on location-consistency rather than scale-consistency. Formal details---including the $L^2$ minimisation proof, the stochastic-observation consistency definition, numerical verification, and the full scale-consistency analysis---are in Appendix~\ref{sec:appendix_consistency}.
